\section{Related Work}

The security of Large Language Models (LLMs) and their agentic counterparts has seen burgeoning interest. Prior work has laid foundational frameworks for understanding and benchmarking these vulnerabilities. \textit{InjecAgent} \cite{zhan2024injecagent} provided early methodologies for benchmarking indirect prompt injections within tool-integrated LLM agents, establishing that indirect vectors pose a severe risk when agents autonomously utilize external tools. The \textit{AgentDojo} \cite{debenedetti2024agentdojo} family of work expanded this by providing a dynamic environment to evaluate both attacks and defensive mechanisms for LLM agents, becoming a standard for indirect prompt injection benchmarking. Furthermore, the \textit{HackAPrompt} competition \cite{schulhoff2023ignore} exposed systemic vulnerabilities in LLMs to diverse prompt hacking techniques.

In the domain of code generation, Pearce et al. \cite{pearce2021asleep} performed an early assessment of GitHub Copilot's security, identifying a propensity for models to reproduce vulnerable coding patterns. This was further formalized by the \textit{Purple Llama CyberSecEval} \cite{bhatt2023purple}, which introduced a secure coding benchmark for LLMs. A 2025 industry report by GitGuardian \cite{gitguardian2025copilot} highlighted a measurable increase in secret leakage rates within Copilot-assisted repositories (6.4\% compared to a 4.6\% baseline). Moreover, the threat of dependency confusion in LLM-generated code echoes the fundamental software supply chain vulnerabilities disclosed by Birsan in 2021 \cite{birsan2021dependency}.

Most recently, \textit{SecRepoBench} \cite{secrepobench2025} benchmarked code agents for secure code completion within real-world repositories, serving as the closest existing analog to our approach. However, as noted in \textit{The Balkanization of Execution-Security Research for AI Coding Agents} \cite{balkanization2026}, existing research remains fragmented. Our study bridges this gap by providing a cross-agent comparison (evaluating Chat, CLI, and IDE environments) across four distinct vulnerability classes operating on a single, shared, realistic codebase.
